% Shared semantic body used by both standard and margin profile wrappers.
% No documentclass and no profile-specific visual code is allowed here.

\begin{titlepage}
  \centering
  \vspace*{0.18\textheight}
  {\Huge\bfseries\sffamily I.P.A.R.A Handbook\par}
  \vspace{0.7em}
  {\Large Shared Semantic Body Regression\par}
  \vspace{1.5em}
  {\small Same content source, different profile renderer\par}
\end{titlepage}
\tableofcontents

\section{对象、表示与不变量}

\subsection{MainText 必须先拥有逻辑主链}

设线性映射 $T\colon V\to W$。选定基以后得到矩阵 $A$，但矩阵只是表示；换基可以改变 $A$，并不改变映射本身。这个最小例子同时检查中文、English、数学字体、section mark 和正文节奏。

\semanticmargin{Owner}{MainText owns the reasoning chain. Margin only carries navigation, reminders and local references.}

English text should remain calm inside a Chinese technical handbook. Inline code such as \code{profile=standard} or \code{profile=margin} should be visible but subordinate to the reasoning.

\begin{mentalmodel}{Mental Model · 表示不是对象}
真正稳定的理解不是记住某个矩阵长什么样，而是先问：哪些量依赖表示，哪些结构在合法变换下保持不变？
\end{mentalmodel}

\begin{mechanism}{Mechanism · 从对象到表示再回到不变量}
先固定数学对象，再选择表示；所有计算都发生在表示层，结论最后必须翻译回对象层。这样可以避免把“坐标发生变化”误读成“对象发生变化”。
\end{mechanism}

\begin{processchain}
  \processstage{Target}
  \processstage{Auxiliary Function}
  \processstage{Qualification}
  \processstage{Boundary / Zero Design}
  \processstage{Witness}
  \processstage{Translate Back}
\end{processchain}

\semanticmargin{Boundary}{表示变化 $\neq$ 对象变化。这个提醒可以进入边栏，但对象/表示的第一次完整定义必须留在 MainText。}

\begin{criterion}{Criterion · 什么时候可以把两种表示视为同一对象}
必须能指出连接两种表示的合法变换，并说明题目真正关心的性质在该变换下保持不变。
\end{criterion}

\begin{boundary}{Boundary · 相似 $\neq$ 合同}
相似变换服务于同一线性算子的不同基表示；合同变换服务于二次型在变量替换下的表示。两个关系都涉及可逆矩阵，但保留的不变量不同。
\end{boundary}

\begin{warning}{Warning · 不要从符号长相推断 Owner}
看到两个式子都出现 $P^{-1}AP$ 或 $P^{\mathsf T}AP$ 时，先识别变换语义，再决定应该调用谱不变量还是惯性不变量。
\end{warning}

\begin{example}{Example · 一个最小运行实例}
若 $B=P^{-1}AP$，则 $A$ 与 $B$ 有相同的特征多项式：
\[
  \det(\lambda I-B)
  =\det\!\left(P^{-1}(\lambda I-A)P\right)
  =\det(\lambda I-A).
\]
\end{example}

\subsection{公式、表格与代码}

较长公式用来观察正文宽度和 display rhythm：
\begin{align}
  \operatorname{Cov}(X,Y)
  &=\mathbb E\!\left[(X-\mathbb EX)(Y-\mathbb EY)\right] \\
  &=\mathbb E(XY)-\mathbb E(X)\mathbb E(Y).
\end{align}

\begin{handbooktable}{L{0.20\linewidth}YY}
\toprule
\textbf{层} & \textbf{拥有的东西} & \textbf{不应该拥有的东西}\\
\midrule
KOMA & section hierarchy、页眉页脚、TOC、duplex mechanics & Mental Model、Boundary 等学习语义\\
CTeX & 中文语言、CJK 与中文命名基础设施 & Handbook margin 页面状态\\
I.P.A.R.A & Semantic API、Profile 与领域组件 & 重写一套底层 book engine\\
\bottomrule
\end{handbooktable}

% Same source: standard keeps ordinary width, margin temporarily consumes outer envelope.
\begin{widecontent}
\begin{boundarytable}
\boundarytablehead
对象 & $\neq$ & 表示 & 对象回答“是什么”；表示回答“怎样计算/编码” & 把坐标变化误当成对象变化\\
相似 & $\neq$ & 合同 & 前者保留线性算子谱结构；后者保留二次型惯性 & 调错不变量与标准形\\
\bottomrule
\end{boundarytable}
\end{widecontent}

\begin{handbooklongtable}{L{0.22\linewidth}YY}
\toprule
\textbf{跨页表职责} & \textbf{推荐做法} & \textbf{避免做法}\\
\midrule
短标签 + 长解释 & 固定短标签列，解释列使用 Y 自动分享剩余宽度 & 让多个固定列宽直接吃满整张表，却忘记列间 padding\\ 
可增长内容 & 允许 xltabular 在需要时分页 & 为了“永不分页”整体缩小字号\\
\bottomrule
\end{handbooklongtable}

\begin{lstlisting}
profile = "standard | margin"
owner = "common/latex/README.md"
renderer = "semantic, not color-named"
\end{lstlisting}

\subsection{TikZ 基础词汇}

\begin{handbookdiagram}
\begin{tikzpicture}[node distance=12mm and 14mm]
  \node[ipara process] (object) {Object};
  \node[ipara process,right=of object] (repr) {Representation};
  \node[ipara process,below=of $(object)!0.5!(repr)$] (inv) {Invariant / Boundary};
  \draw[ipara edge] (object) -- node[ipara label,above] {choose} (repr);
  \draw[ipara edge] (repr) -- node[ipara label,right] {audit} (inv);
  \draw[ipara ref] (inv) -- (object);
\end{tikzpicture}
\diagramcaption{Core TikZ vocabulary regression}
\end{handbookdiagram}

\section{分页、目录与引用}

这一节故意存在，用于检查 TOC、section/subsection marks 与跨页稳定性。

\subsection{验收边界}

\semanticmargin{Stop Boundary}{Supplementary channel can move between inline and outer margin, but it cannot become the only owner of a necessary premise.}

\begin{enumerate}
  \item 不出现 CTeX non-standard heading warning；
  \item 不出现 `fancyhdr`，页眉页脚只由 `scrlayer-scrpage` 拥有；
  \item portable fonts 全部解析；
  \item semantic names 与 renderer 分离；
  \item `semanticmargin` 在 standard 中 inline、在 margin 中进入 outer margin；
  \item `widecontent` 在 standard 中安全退化、在 margin 中扩展到 outer envelope；
  \item 没有 unexplained overfull/underfull、undefined reference 或 package conflict。
\end{enumerate}
